% 本文件根据开题报告与中期进展报告整理，可直接作为毕业论文绪论章节引用。

\chapter{绪论}
\label{chap:introduction}

\section{研究背景与意义}

随着``工业4.0''、智能制造和工业互联网的持续推进，现代工业系统正朝着大型化、复杂化、智能化方向快速发展。电源系统、逆变器系统以及典型流程工业装置中集成了大量传感器、执行器与控制单元，系统内部同时存在电磁、热力、机械和控制等多重耦合关系，运行过程具有显著的动态性、非线性和不确定性。一旦关键设备退化、控制偏差累积或外部工况突变未能被及时识别，就可能诱发局部失稳、连锁故障甚至停机事故，进而造成设备损坏、产品质量下降和严重的经济损失。因此，面向复杂工业过程开展高可靠异常检测研究，具有重要的理论价值和工程意义。

在电源系统及其相关工业场景中，异常往往表现为电压、电流、功率、温度等多类变量的协同变化。与一般静态监测任务不同，这类系统中的异常既可能体现为瞬态尖峰、缓慢漂移，也可能表现为谐波畸变、频率偏移和波形失真等复杂模式，单纯依靠人工巡检或固定阈值规则已难以满足实时监测和早期预警需求。近年来，数据驱动方法依托工业现场不断积累的运行数据，为异常检测提供了新的研究路径。通过挖掘多变量时间序列中的潜在演化规律与变量关联结构，可以在缺乏完备机理模型的条件下实现对设备健康状态和过程异常的有效识别。

然而，真实工业数据通常并非理想的单采样率序列。受传感器成本、测量原理、通信带宽以及变量动态特性差异等因素影响，不同变量常以不同频率采集。例如，电流、电压等电气量往往以较高频率采样，而温度、状态量或质量指标则通常以较低频率更新。将这些数据统一到同一时间轴后，会产生大量规则性缺失值和时间异步现象，形成典型的多采样率时间序列。多采样率特性不仅削弱了数据完整性，也显著增加了建模难度：如果直接套用单采样率异常检测方法，往往难以充分利用低频变量信息，甚至会因缺失值传播、插补失真和时序错位而降低检测精度。

基于上述背景，围绕多采样率工业过程数据开展异常检测研究，特别是探索``数据插补与异常检测协同优化''的建模思路，具有双重意义。从理论层面看，该研究有助于深化对多采样率时间序列表示、缺失信息恢复和异常模式判别机制的认识，丰富工业过程监测与智能诊断的研究体系。从工程层面看，通过软件建模手段提升对低频变量和隐性异常的感知能力，可在一定程度上降低对高成本高频传感器的依赖，为电源系统及相关工业场景的安全运行、故障预警和运维决策提供支持。

\section{国内外研究现状}

\subsection{传统异常检测方法}

异常检测早期研究主要基于统计学建模、距离与密度度量、聚类分析以及支持向量机等传统方法。统计学方法通过假设数据服从某类分布并设置判别阈值实现异常识别，具有实现简单、计算开销较低等优点；距离、密度和聚类类方法则通过样本间相似性关系识别离群行为；单类支持向量机等判别式方法在缺乏异常样本时也具有一定应用价值。这类方法在结构较简单、维度较低和分布相对稳定的场景中取得了较好效果。

但对于复杂工业过程而言，传统方法的局限性也较为明显。首先，多数方法依赖显式分布假设或人工设计特征，当系统呈现强非线性、强耦合和时变特性时，模型鲁棒性往往不足。其次，在高维多变量场景下，距离和密度的判别能力会受到``维度灾难''影响，聚类效果对参数和特征选择也较为敏感。再次，传统方法通常难以刻画长时间跨度内的动态依赖关系，对复杂异常模式和多类型故障的适应能力有限。因此，随着工业数据规模和复杂度不断提升，仅依靠传统方法已难以满足高精度、强泛化和实时异常检测的需求。

\subsection{基于深度学习的异常检测方法}

深度学习方法能够通过端到端训练自动提取时序特征，在高维、非线性和强噪声场景中展现出更强的建模能力，已成为异常检测领域的重要研究方向。现有研究大体可以分为两类：一类是基于预测的方法，即利用历史序列预测未来观测值，并将预测误差作为异常判据；另一类是基于重构的方法，即通过自编码器、变分自编码器或 Transformer 等模型学习正常样本的潜在表示，再根据重构误差识别异常。

以 LSTM、TCN 和 Transformer 为代表的深度时序模型在工业异常检测中得到了广泛应用。LSTM 擅长建模时间依赖关系，TCN 具有并行计算效率高、感受野可扩展等特点，Transformer 则依靠自注意力机制在长序列建模方面表现突出。与传统方法相比，深度学习模型减轻了对人工特征工程的依赖，更适合处理复杂工业系统中的高维动态数据。

尽管如此，现有深度学习异常检测方法仍然存在若干不足。其一，多数方法默认所有变量具有统一采样率，对多采样率场景下的规则性缺失和时间异步问题考虑不足。其二，很多模型主要聚焦时间维度建模，对变量间耦合关系、系统拓扑结构以及频域信息的利用仍不充分。其三，在实际部署中，异常阈值的设定、噪声干扰的抑制以及模型可解释性等问题仍然是影响检测性能的重要因素。

\subsection{多采样率过程建模与异常检测研究}

针对工业现场广泛存在的多采样率问题，近年来已有研究开始探索适配性更强的建模方法。基于因子分析的多采样率建模方法能够从统计相关性角度提取潜在变量，对多采样率过程故障检测提供了较早的研究思路；层次化循环神经网络方法如 MRA-LSTM 通过多尺度状态更新机制建模不同采样频率下的时序依赖，增强了对高频变量和低频变量协同关系的表达能力；基于 Transformer 的多采样率模型则通过数据重组、采样类型编码以及预训练--微调策略，提高了对复杂时序关系和缺失模式的建模效果。

总体来看，多采样率异常检测研究已经从简单的数据对齐与统计分析，逐步发展到面向深度特征学习的统一建模阶段。但是，现有方法仍然面临一些共性挑战：一是低频变量长期缺失容易导致插补结果过度平滑或偏离真实物理状态；二是变量间的耦合关系往往隐式存在且可能随工况变化而变化，固定结构难以充分反映真实关联；三是电源系统等场景中的周期性和谐波异常在频域中更具可分性，单纯依赖时域建模容易漏检；四是传统``先插补、后检测''的两阶段流程容易出现目标不一致问题，即插补结果在数值上更平滑，却削弱了对异常特征的敏感性。上述问题表明，多采样率异常检测仍有进一步研究和改进的空间。

\section{现有研究存在的主要问题}

结合开题阶段的文献调研与中期阶段的模型实现，可以将当前多采样率异常检测研究中的关键难点概括为以下几个方面：

\begin{itemize}
    \item \textbf{规则性缺失与时间异步并存。} 多采样率数据中的缺失并非随机噪声，而是由采样频率差异带来的结构性缺失。低频变量在统一时间轴上长期缺测，容易造成序列信息不完整、监督信号稀疏以及样本对齐困难。
    \item \textbf{变量耦合关系复杂且拓扑未知。} 电源系统及工业过程中的变量之间往往存在显著的物理耦合和控制耦合，但真实拓扑结构难以完整获取，固定邻接关系或独立建模策略都难以兼顾表达能力与泛化性能。
    \item \textbf{周期性异常和频谱变化难以在纯时域中显式刻画。} 对于谐波畸变、频率漂移、波形失真等异常，频域特征往往比时域波形更敏感。如果模型只在时域进行插补和检测，容易忽略关键频谱信息。
    \item \textbf{插补目标与检测目标存在偏差。} 传统两阶段流程通常将插补视为独立预处理步骤，优化重点是恢复数值本身，而非保留异常判别所需的结构信息，因而可能在无意中削弱异常特征。
    \item \textbf{异常评分与阈值机制鲁棒性不足。} 在无监督或弱监督场景下，异常阈值常依赖经验设定或静态统计量，当工况发生变化时容易出现误报和漏报，影响模型的实际部署效果。
\end{itemize}

\section{本文的研究内容与技术路线}

针对上述问题，本文面向电源系统等具有显著多采样率特征的工业场景，研究基于数据插补的多变量时间序列异常检测方法，并提出一种自适应图--频率融合异常检测网络 AGF-ADNet（Adaptive Graph-Frequency Anomaly Detection Network）。该方法以``插补服务于检测''为核心思想，将缺失值恢复、变量关系建模、时频特征提取和异常判定统一到同一框架中，以提高多采样率场景下的异常检测性能。

本文的主要研究内容如下：

\begin{itemize}
    \item \textbf{多采样率数据对齐与掩码建模。} 将不同采样频率下的多变量序列统一到共享时间轴上，通过构建缺失掩码描述低频变量的规则性缺失模式，并为后续插补和误差计算提供约束信息。
    \item \textbf{自适应图结构学习。} 利用可学习的节点表示和数据驱动的邻接矩阵生成机制，自动挖掘变量之间的潜在耦合关系，在拓扑未知条件下增强跨变量信息传播能力。
    \item \textbf{时频协同插补。} 在时域分支中结合图卷积网络与多尺度时序卷积网络，刻画变量间依赖和多尺度动态模式；在频域分支中引入傅里叶变换与注意力机制，强化对周期结构和关键频段的表达能力；随后通过门控融合模块实现时域与频域信息的自适应整合。
    \item \textbf{基于重构误差的异常检测。} 将插补后的完整序列输入 Transformer 重构器，利用重构误差构建异常得分，并结合统计阈值策略实现异常判定，从而提升对复杂异常模式的识别能力。
    \item \textbf{实验验证与方法比较。} 以典型多采样率工业过程数据为对象，对所提方法与 LSTM、MRA-LSTM、Multirate-Former 等代表性方法进行比较，采用准确率、精确率、召回率和 F1-Score 等指标评估模型性能。
\end{itemize}

本文的整体技术路线可以概括为：首先对多采样率数据进行标准化与滑动窗口构造；其次在统一时间轴上结合掩码信息完成多变量序列建模；随后通过自适应图学习和时频双分支插补恢复缺失信息，并利用 Transformer 对完整序列进行全局重构；最后根据重构误差计算异常得分，完成异常识别与效果评估。该技术路线兼顾了多采样率数据的缺失特征、变量耦合特征与频谱结构特征，能够较系统地解决传统方法在复杂工业场景下面临的建模困难。

\section{本文的主要创新点}

围绕多采样率工业过程异常检测中的关键问题，本文的创新点主要体现在以下几个方面：

\begin{itemize}
    \item \textbf{构建了面向多采样率场景的插补--检测一体化框架。} 不再将数据插补仅视为独立预处理步骤，而是通过统一模型结构使插补结果直接服务于异常判定，提高整体建模的一致性与鲁棒性。
    \item \textbf{提出了基于自适应图学习的变量耦合建模机制。} 在缺乏先验拓扑信息的条件下，模型能够自动发现变量间的潜在关联关系，为低频变量恢复和异常识别提供结构约束。
    \item \textbf{引入了频域增强的时频协同建模思路。} 通过在插补过程中融合频域注意力机制，增强模型对周期性、谐波性异常模式的感知能力，为电源系统等具有显著频谱特征的场景提供了更有针对性的解决方案。
\end{itemize}

\section{论文结构安排}

除本章外，全文其余部分拟安排如下：

\begin{itemize}
    \item 第二章介绍多采样率时间序列建模、图神经网络、时序卷积、傅里叶变换与 Transformer 等相关理论与关键技术，为后续方法设计奠定基础。
    \item 第三章详细阐述面向多采样率异常检测的 AGF-ADNet 模型，包括数据表示方式、图结构学习机制、时频双分支插补策略以及异常评分方法。
    \item 第四章给出实验设计、评价指标与对比方法，并对模型在典型多采样率工业过程数据上的检测结果进行分析，验证所提方法的有效性。
    \item 第五章对全文工作进行总结，归纳研究结论，并结合当前方法的局限性展望后续改进方向。
\end{itemize}
